\begin{solapartat}
\punts{0,25} De la gràfica, s'han completat 7 voltes.

\punts{1,0} De l'esquema la longitud de l'eix major és: $35R_S + 2R_S + 200R_S = 237R_S$

Si el resultat és $235R_S$ perquè no s'ha tingut en compte el diàmetre del Sol, llavors la qualificació de la resposta serà 0,8~p en lloc d'1~p.
\figura[0.6]{sol_fig1.png}
\end{solapartat}

\begin{solapartat}
\punts{0,25} Dibuix d'una òrbita el·líptica amb el sol en un dels seu focus.

\punts{0,1} Identificar els punts A i C com el més proper i el més allunyat del Sol.

\punts{0,25} Representar correctament la direcció i sentit dels vectors acceleració i velocitat.

\medskip
\punts{0,4} Conservació de l'energia mecànica:
% DUBTE: la pauta escriu M_T (massa de la Terra) tot i que la massa central és la del Sol.
\[ E_m = \frac{1}{2}mv^2 - G\,\frac{mM_T}{r} \Longrightarrow \frac{1}{2}mv^2 = \mathit{constant} + G\,\frac{mM_T}{r} \]

\punts{0,25} Quan $r$ és mínim, llavors el terme $G\,\frac{mM_T}{r}$ i l'energia cinètica són màxims, per tant, el punt on la velocitat és més gran és el punt més proper al Sol, el punt A. També es pot argumentar a partir de la conservació de l'energia mecànica que l'energia cinètica és màxima quan l'energia potencial gravitatòria és mínima, i això passa quan la distància entre el Sol i la nau és mínima.

\medskip
\destaca{Alternativament}

\punts{0,4} Segona llei de Kepler: ``Un segment rectilini que uneix la nau i el Sol escombra àrees iguals en intervals de temps iguals''.

\punts{0,25} Com que la longitud d'aquest segment és mínima al punt A, llavors, per escombrar la mateixa àrea en el mateix interval de temps, cal que el desplaçament sigui màxim, és a dir, cal que la velocitat de la nau sigui màxima.
\end{solapartat}
